\section{The symbolic engine}\label{sec:symbolic}
\begingroup\sloppy

\providecommand{\Reach}{\mathit{Reach}}
\providecommand{\Uni}{U}
\providecommand{\pre}{\mathrm{pre}}
\providecommand{\post}{\mathrm{post}}
\providecommand{\Bad}{\mathit{Bad}}
\providecommand{\Ok}{\mathit{Ok}}
\providecommand{\lddsem}[1]{[\![#1]\!]}

This section describes the engine that decides every goal:
\code{strong_sim}, \code{strong_bisimilar}, \code{weak_sim} and
\code{weak_bisimilar} between closed start terms, without building an
explicit graph of the systems. It reads the user's LTS as a set of rules,
splits each state into a finite \emph{shape} and a vector of \emph{slot}
values, learns one transition relation per \emph{summand} (a rule chain) on
only the slots it reads, stores sets and relations as
Sylvan's~\cite{vandijk-vdpol15} List Decision Diagrams
(LDDs)~\cite{blom-vdpol08}, and decides the weak or the strong game as a
greatest fixpoint over pairs. With the log on, the game's last check is
repeated with Sylvan's proof log open: that log is the certificate
(\cref{sec:prooflog}). The engine is not trusted; the kernel checks the
certificate.

\paragraph{How a user reaches it.}
\begin{itemize}
\item The tactics \code{sim}, \code{bisim} (strong) and \code{weak_sim},
  \code{weak_bisim} (weak) (\file{src/g_mebi.mlg}) all call
  \code{Learn.check ~strong:(not weak)} (\file{src/user_tactics.ml},
  \code{read_goal}). With the prefix \code{debug} they also print
  \code{Learn.pr} (\cref{sec:sym-debug}).
\item \code{Print MeBi Rules <lts>} and
  \code{Print MeBi Skeleton <term> Using <lts>} (\file{src/g_mebi.mlg})
  print the first two stages.
\end{itemize}

\subsection{Map of the code}

\begin{center}
\small
\begin{tabularx}{\linewidth}{@{}l r >{\raggedright\arraybackslash}X@{}}
\toprule
File & lines & Role \\
\midrule
\file{lib/rules/rules.ml}, \file{.mli} & 209, 73 &
  An LTS's constructors read as rules (\cref{sec:sym-rules}). \\
\file{lib/skeleton/skeleton.ml}, \file{.mli} & 786, 138 &
  Structural constructors, shapes and slots, summands
  (\cref{sec:sym-skeleton}). \\
\file{lib/ldd/ldd.ml}, \file{.mli} & 123, 96 &
  OCaml interface to Sylvan's LDDs (\cref{sec:sym-ldd}). \\
\file{lib/ldd/ldd_stubs.c} & 235 & The C stubs. \\
\file{lib/ldd/dune} & 48 & Compiles Sylvan and Lace into the binding's
  library. \\
\file{vendor/sylvan/}, \file{vendor/lace/} & (C) & git submodules
  (forks): Sylvan on the branch \code{proof-log} (the proof log), Lace
  1.10.0 (\file{lib/ldd/README.md}). \\
\file{lib/learn/learn.ml}, \file{.mli} & 1195, 79 &
  The interpreter, learning, exploration, the game, the logged last check
  (\cref{sec:sym-learn,sec:sym-game,sec:sym-check}). \\
\file{test/ldd_tests.ml} & & 25 checks of the binding and the proof
  log, no Rocq; run in CI (\cref{sec:tests}). \\
\file{src/user_tactics.ml} & & \code{prove}: the four tactics, the call to
  \code{Learn.check}, the log's place; \code{check_log}. \\
\bottomrule
\end{tabularx}
\end{center}

The libraries depend on each other in a line:
\code{rocq-mebi.rules} $\leftarrow$ \code{rocq-mebi.skeleton}
$\leftarrow$ \code{rocq-mebi.learn} $\rightarrow$ \code{rocq-mebi.ldd}
(\file{lib/learn/dune}, \file{lib/skeleton/dune}). Only \code{ldd} is free
of \code{rocq-runtime}, which is why its test binary needs no Rocq.

\begin{figure}[t]
\centering
\begin{tikzpicture}[node distance=6mm and 7mm]
  \node[box] (lts) {Rocq LTS\\ \scriptsize inductive $M\to\Opt\,A\to M\to\mathsf{Prop}$};
  \node[box,right=of lts] (rules) {rules\\ \scriptsize \texttt{Rules.of\_lts}};
  \node[box,right=of rules] (skel) {shapes, slots,\\ summands\\ \scriptsize \texttt{Skeleton}};
  \node[box,right=of skel] (enc) {vectors\\ $[\mathit{shape};v_0;\dots;v_{w-1}]$\\ \scriptsize \texttt{encode}};
  \node[box,below=12mm of enc] (learn) {learning\\ \scriptsize \texttt{fire}, \texttt{interp}};
  \node[box,left=of learn] (rel) {relations per\\ summand, label,\\ meta (LDDs)};
  \node[box,left=of rel] (expl) {chained\\ exploration\\ \scriptsize \texttt{relprod}};
  \node[box,left=of expl] (game) {weak or strong game\\ $\nu F$ on $\Reach_A\times\Reach_B$\\ \scriptsize \texttt{relprev}};
  \node[box,below=10mm of game] (verdict) {verdict};
  \node[box,right=of verdict] (cert) {proof log\\ \scriptsize (Sylvan)};
  \node[box,right=of cert] (coq) {facts, typechecked\\ \scriptsize \cref{sec:prooflog}};
  \draw[dep] (lts) -- (rules);
  \draw[dep] (rules) -- (skel);
  \draw[dep] (skel) -- (enc);
  \draw[dep] (enc) -- (learn);
  \draw[dep] (learn) -- (rel);
  \draw[dep] (rel.170) -- (expl.10);
  \draw[dep] (expl.-10) -- node[lbl,below]{frontier} (rel.-170);
  \draw[dep] (expl) -- (game);
  \draw[dep] (game) -- (verdict);
  \draw[dep] (verdict) -- (cert);
  \draw[dep] (cert) -- (coq);
\end{tikzpicture}
\caption{The symbolic pipeline, and the certificate it leaves.}
\label{fig:sym-pipeline}
\end{figure}

\subsection{Rules: \texttt{lib/rules}}\label{sec:sym-rules}

\begin{definition}[Rule]
Let \code{L : M -> option A -> M -> Prop} be an inductive without
parameters (checked by \code{Rocq_utils.lts_shaped}). Each constructor
\[
  c : \forall\, x_1\dots x_k,\ P_1 \to \dots \to P_m \to L\ s\ \ell\ t
\]
is a rule $\rho_c$ with variables $\vec x$ (the binders that are not
propositions), premises $P_1,\dots,P_m$ (the binders that are), source
$s$, label $\ell$ and target $t$. A premise is
\begin{itemize}
\item \code{Step}: a step $L'\ s'\ \ell'\ t'$ of an LTS $L'$ (any LTS, not
  necessarily $L$);
\item \code{Decide}: a proposition $P$ for which typeclass resolution finds
  an instance of stdlib's \code{Decidable P}; the rule keeps
  \code{Decidable_witness P i}, a boolean over the rule's variables;
\item \code{Opaque}: any other proposition (an equation, \code{In}, a
  relation, an existential).
\end{itemize}
\end{definition}

The types are at \locr{lib/rules/rules.mli}{15}{51}. Every term of a rule
is over the \emph{whole} context \code{ctx} in de Bruijn indices: each
premise, read in the context of the binders before it, is lifted over the
rest (\loc{lib/rules/rules.ml}{108}). \code{rule}
(\locr{lib/rules/rules.ml}{87}{139}) classifies each binder: a non-\code{Prop}
binder is a variable (\loc{lib/rules/rules.ml}{110}); a \code{Prop} binder
is a \code{Step} if its head is an LTS-shaped inductive applied to three
arguments (\code{step_args}, \locr{lib/rules/rules.ml}{68}{78}), else a
\code{Decide} if \code{witness} (\locr{lib/rules/rules.ml}{50}{63}) finds a
\code{Decidable} instance, else \code{Opaque}
(\loc{lib/rules/rules.ml}{122}). The conclusion must be a step of the
same inductive, or \code{Not_an_lts} is raised. \code{of_lts}
(\locr{lib/rules/rules.ml}{142}{166}) returns the rules of the root LTS
and, transitively, of every LTS its \code{Step} premises use, first met
first. The \code{Decidable} class must be loaded:
\file{theories/loader.v} \code{Require}s (does not import)
\code{Stdlib.Classes.DecidableClass}.

\code{Opaque} premises are kept; the interpreter
stops at them (\cref{sec:sym-interp}) and the skeleton marks the summand
dynamic.

\begin{example}[\code{Print MeBi Rules}]
\begin{rocq}
Inductive step : tm -> option act -> tm -> Prop :=
| s_pre x k : step (pre x k) (Some x) k
| s_hide k l k' : step k l k' -> visible l = true -> step (hide k) l (hide k').
\end{rocq}
reads as two rules; \code{s_hide} has a \code{Step} premise and a
\code{Decide} premise (\code{=} on \code{bool} has a stdlib instance). The
rule \code{q : forall n m, (exists k, k = n) -> opaque n None m} has an
\code{Opaque} premise, printed with \code{[Opaque]}
(\loc{lib/rules/rules.ml}{180}).
\end{example}

\paragraph{Stale comments.} \file{lib/rules/dune} says ``only the plugin
(src/) uses it'' and ``S3's bit format goes in a library of its own'';
\file{doc/s2-linearisation.md} decision 8 says the same. Today
\file{lib/skeleton}, \file{lib/learn} and \file{lib/certify} use it, and
the bit format is gone.

\subsection{The skeleton: \texttt{lib/skeleton}}\label{sec:sym-skeleton}

Phase~1 turns a start term into a vector whose positions a rule touches
locally. Everything is computed from the user's rules, with no knowledge
of any calculus (\file{doc/skeleton.md}).

\subsubsection{Structural constructors}

\begin{definition}[Structural]
For a rule $\rho$ with source $s$ and target $t$, a constructor $c$
occurring in $s$ at path $\pi$ (a list of field indices, parameters
excluded) is \emph{kept} by $\rho$ if $t$ has a constructor application
with head $c$ at the same path $\pi$. $c$ is \emph{structural} for $L$ if
some rule of $L$ keeps it. Constructors that occur in sources but that no
rule keeps are \emph{never kept}.
\end{definition}

\code{classify} (\locr{lib/skeleton/skeleton.ml}{103}{125}) computes both
lists from \code{occurrences} and \code{at}
(\locr{lib/skeleton/skeleton.ml}{81}{96}). Only the rules of the root LTS
are used (\code{structural_of}, \locr{lib/skeleton/skeleton.ml}{726}{733}).
CCS's \code{par} is structural (kept by \code{s_parl}); Proc's \code{cpar}
is structural although \code{do_par_end} dissolves it, since
\code{do_parl} keeps it. \code{pre}, \code{tend} are never kept.

An earlier version (same day) took the constructors that \emph{every} rule
keeps and collapsed a node into one slot where a rule changed its shape.
Proc then had no skeleton at all, so the shape became part of the state
instead (decision D3 of \file{SYMBOLIC_PLAN.md}).

\subsubsection{Shapes, slots and the canonical split}

\begin{definition}[Shape, split]
A \emph{shape} is a tree $\sigma ::= \mathsf{Slot}\ i \mid
\mathsf{Node}(c, \vec p, \sigma_1\dots\sigma_n)$ with $c$ structural and
$\vec p$ its type parameters (\locr{lib/skeleton/skeleton.ml}{5}{14}).
The \emph{canonical split} of a closed term $t$ is
\[
\mathrm{split}(t) =
\begin{cases}
\mathsf{Node}(c,\vec p,\mathrm{split}(t_1)\dots\mathrm{split}(t_n)) &
  t = c\ \vec p\ t_1\dots t_n \text{ fully applied},\ c \text{ structural}\\
\mathsf{Slot}\ i,\ \text{with value } t & \text{otherwise}
\end{cases}
\]
slots numbered $0,1,\dots$ left to right (\code{split},
\locr{lib/skeleton/skeleton.ml}{736}{753}; inverse \code{rebuild},
\locr{lib/skeleton/skeleton.ml}{756}{765}). A slot's value never starts
with a structural constructor. Shapes are compared by \code{key}
(\locr{lib/skeleton/skeleton.ml}{134}{142}), a string of the structure
with slots unnumbered.
\end{definition}

A \emph{state} is a shape together with its slot values. Since a slot
value is never structural, a rule's structural constructors match nodes
only, and every other part of its source reads slot values.

\paragraph{A possible key collision.} \code{key} prints a constructor as
its mutual block's name and the constructor's index
(\code{MutInd.to_string (fst (fst n.cons))}, \code{snd n.cons}); it omits
the index of the inductive within the block, and the parameters. Two
structural constructors with the same index in two inductives of one
mutual block, or one polymorphic constructor at two parameter
instantiations, would get the same key. No example in the repository has
either; it is a latent issue, not an observed one.

\subsubsection{Summands}

\begin{definition}[Summand, origins]
A \emph{summand} of a shape $\sigma$ (\locr{lib/skeleton/skeleton.ml}{31}{37})
is a rule chain that can fire from $\sigma$: a list of rule names,
outermost first, where \code{<L>} stands for ``a step of a closed value by
$L$, computed at run time''; the source slots it \emph{reads}; its
\emph{target} $\tau ::= \mathsf{TSlot}(o) \mid \mathsf{TNode}(c,\vec
p,\tau_1\dots\tau_n)$ with origins $o ::= \mathsf{Moved}\ i \mid
\mathsf{Computed}\ \vec \imath$ (the slot's value is source slot $i$
unchanged, or computed from slots $\vec\imath$); the label's head
constructor where the rules fix it; and a flag \code{dynamic}, set when a
computed value may start with a structural constructor at run time, so
its shape is known only then.
\end{definition}

\code{summands} (\locr{lib/skeleton/skeleton.ml}{381}{535}) computes them
by abstract matching of each rule against the shape, seen as a target of
moved slots (\code{target_of_shape}):

\begin{shell}[mathescape=true]
go($\tau$):                                            (skeleton.ml:394-427)
  if $\tau$ = TSlot o:  [ chain <L>, reads(o), target = normalize(Computed reads(o)) ]
  if $\tau$ is a node already in progress:  [ chain <L>, Computed, dynamic ]   (recursion guard)
  else  concat over rules $\rho$ of L: rule($\rho$, $\tau$)
rule($\rho$, $\tau$):                                       (skeleton.ml:428-533)
  pmatch $\rho$.src against $\tau$          -- No_match: structural c vs other node/slot
  for the premises, in order:
    Decide P / Opaque P:  reads += slots P mentions   (Opaque: dynamic)
    Step L' s' l' t':
      if L' = L and inst(s') is made of nodes and moved slots:
        for each summand $\mu$ of go(inst(s')), compatible label head:
          pmatch t' against $\mu$.target; chain @= $\mu$.chain; reads += $\mu$.reads
      else  chain @= [<L'>]; t' bound to "computed from inst(s')'s slots"
  target := normalize(inst($\rho$.tgt));  reads += label's slots + computed_reads(target)
\end{shell}

\begin{itemize}
\item \code{pmatch} (\locr{lib/skeleton/skeleton.ml}{224}{246}) binds a
  variable to a part of the shape; a variable bound twice (a non-linear
  pattern) reads both parts (equality decided at run time); a
  non-structural constructor over a slot reads the slot.
\item \code{inst} (\locr{lib/skeleton/skeleton.ml}{251}{277}) turns a rule
  term into a target: structural constructors become nodes, a variable
  bound to a part of the shape is moved, anything else is computed from
  the slots its variables are bound to.
\item \code{normalize} (\locr{lib/skeleton/skeleton.ml}{335}{365}) expands
  a computed value whose type has a single, structural constructor (a
  pair, a record) into its node, with fuel 20, and reports whether a
  computed value's type has structural constructors (\code{carries},
  \locr{lib/skeleton/skeleton.ml}{321}{328}): then the summand is
  \emph{dynamic}.
\item \code{computed_reads} (\locr{lib/skeleton/skeleton.ml}{187}{191})
  adds the sources of each computed slot to the reads. It was missing in
  stage 2b, and made CADP's \code{DO_ACQUIRE} get stuck
  (\file{doc/learn.md}, stage 2c, item 2).
\end{itemize}

\code{build} (\locr{lib/skeleton/skeleton.ml}{539}{598}) normalises the
start term with \code{nf_all}, splits it, and explores shapes
breadth-first by following the summands' target shapes, up to
\code{max_shapes} $=200$; it is used only by \code{Print MeBi Skeleton}.
The engine instead asks for summands per shape as it meets shapes
(\code{summands_of}, \locr{lib/skeleton/skeleton.ml}{768}{777}), without
a bound. \code{pr} (\locr{lib/skeleton/skeleton.ml}{621}{721}) prints the
classes, the start shape and its first 12 values, the first 15 shapes,
and the locality: the number of summands, how many keep their shape and
the average number of slots they write (\code{writes},
\locr{lib/skeleton/skeleton.ml}{604}{619}), how many change it, how many
are dynamic.

\begin{example}[\code{Print MeBi Skeleton}]\label{ex:sym-par}
\begin{rocq}
Inductive pr := nil | pre (x : act) (k : pr) | par (p q : pr).
Inductive step : pr -> option act -> pr -> Prop :=
| s_pre x k : step (pre x k) (Some x) k
| s_parl p q l p' : step p l p' -> step (par p q) l (par p' q)
| s_parr p q l q' : step q l q' -> step (par p q) l (par p q')
| s_end : step (par nil nil) None nil.
Print MeBi Skeleton (par (par (pre a nil) (pre b nil)) (pre a nil)) Using step.
\end{rocq}
Derived by hand from the code (not run for this report): \code{par} is
structural, \code{pre} and \code{nil} never kept. The start shape is
$\sigma_1 = \mathsf{par}(\mathsf{par}(P_0,P_1),P_2)$ with
$P_0 = P_2 = \code{pre a nil}$, $P_1 = \code{pre b nil}$. Its summands:
\begin{center}
\small
\begin{tabular}{@{}l l l l@{}}
\toprule
chain & reads & target & keeps shape \\
\midrule
\code{s_parl; s_parl; <step>} & 0 & $\mathsf{par}(\mathsf{par}(C_{0}, M_1), M_2)$ & yes \\
\code{s_parl; s_parr; <step>} & 1 & $\mathsf{par}(\mathsf{par}(M_0, C_{1}), M_2)$ & yes \\
\code{s_parl; s_end} & 0, 1 & $\mathsf{par}(C_{\emptyset}, M_2)$ & no \\
\code{s_parr; <step>} & 2 & $\mathsf{par}(\mathsf{par}(M_0, M_1), C_{2})$ & yes \\
\bottomrule
\end{tabular}
\end{center}
($M_i$ = moved from slot $i$, $C_{\vec\imath}$ = computed from
$\vec\imath$.) The shapes reached are $\sigma_1$, $\sigma_2 =
\mathsf{par}(P_0,P_1)$ (summands \code{s_parl; <step>}, \code{s_parr;
<step>}, \code{s_end}) and $\sigma_3 = P_0$ (one summand \code{<step>}):
3 shapes, 6 summands keeping their shape, 2 changing it, all 8 dynamic,
as \file{doc/skeleton.md}'s table records. Every summand is dynamic
because a computed value of type \code{pr} could be a \code{par}.
\end{example}

\paragraph{Measured locality} (\file{doc/skeleton.md}, 2026-11-06):

\begin{center}
\small
\begin{tabularx}{\linewidth}{@{}>{\raggedright\arraybackslash}p{3.3cm} >{\raggedright\arraybackslash}X r r r r@{}}
\toprule
system & start shape & shapes & keep/ change & avg.\ writes & dyn. \\
\midrule
CCS \code{abp} & \code{res P0 (par (par P1 P2) (par P3 P4))} & 1 & 16/0 & 1.8 & 16 \\
Proc stacked, Test4 \code{p} & \code{cpar (cpar ..) (cpar ..)} & 9 & 39/31 & 1.3 & 0 \\
Proc flat, Test1 \code{p} & \code{P0} (\code{tfix} never kept) & 1 & 1/0 & 1 & 1 \\
CADP fine-grained, 2 processes & \code{pair (PAR ..) ..}, 10 slots & 1 & 2/0 & 6 & 0 \\
CADP fine-grained, 3 processes & 14 slots & 1 & 3/0 & 6 & 0 \\
\code{PrintSkeleton} & \code{par (par P0 P1) P2} & 3 & 6/2 & 1 & 8 \\
streams (\code{nats 0}) & \code{P0} & 1 & 1/0 & 1 & 0 \\
\bottomrule
\end{tabularx}
\end{center}
The CCS row reproduces the ABP's synchronisation vectors (4 local steps,
12 synchronisations); stacked Proc has the 9 binary trees of at most 4
leaves, with \code{do_par_end} and \code{do_assoc} moving between them;
flat Proc has no structure until \code{tfix} unfolds. CADP's ``writes''
over-estimate: every computed field counts as written.

\subsection{List Decision Diagrams: \texttt{lib/ldd}}\label{sec:sym-ldd}

\subsubsection{Definition}

\begin{definition}[LDD; after van Dijk and van de Pol]
An LDD is a node of a directed acyclic graph with two terminals
$\mathbf 0$ and $\mathbf 1$; every other node $v$ has a value
$\mathit{val}(v)\in\N$ (here below $2^{32}-1$), a \emph{down} edge and a
\emph{right} edge. Its meaning, a set of finite sequences over $\N$, is
\[
\lddsem{\mathbf 0} = \emptyset,\qquad
\lddsem{\mathbf 1} = \{\varepsilon\},\qquad
\lddsem{v} = \{\mathit{val}(v)\cdot w \mid w \in \lddsem{\mathit{down}(v)}\}
  \cup \lddsem{\mathit{right}(v)}.
\]
Invariants: $\mathit{down}(v) \ne \mathbf 0$; $\mathit{right}(v)$ is
$\mathbf 0$ or a node with a strictly larger value (so the right chain is
a sorted list of the values possible at one level); nodes are unique in a
hash table. Under these invariants each set has exactly one LDD, so set
equality is node identity.
\end{definition}

In Sylvan, \code{lddmc_makenode(value, ifeq, ifneq)} builds a node with
\code{ifeq} the down edge and \code{ifneq} the right edge; it returns
\code{ifneq} when \code{ifeq} is $\mathbf 0$ and asserts the order
(\file{vendor/sylvan/src/sylvan_ldd.c}, lines 352--364). Values are labels,
never bits: a domain needs no bound and no width. In MeBi each slot value
(a Rocq term of any size or type) gets the next integer from a table, so
the cost of a deep or large term is paid once, in the table, not in the
diagram.

\begin{figure}[t]
\centering
\begin{tikzpicture}[
  n/.style={draw,rectangle,minimum width=6mm,minimum height=5mm,inner sep=1pt,font=\small},
  dn/.style={-{Stealth[length=1.8mm]}},
  rt/.style={-{Stealth[length=1.8mm]},dashed}]
  % level 0: shape
  \node[n] (a1) at (0,0) {1};
  \node[n] (a2) at (2.2,0) {2};
  \node[n] (a3) at (4.6,0) {3};
  % level 1: slot 0
  \node[n] (b1) at (-1.2,-1.3) {1};
  \node[n] (b2) at (0.2,-1.3) {3};
  \node[n] (b3) at (2.2,-1.3) {3};
  \node[n] (b4) at (4.6,-1.3) {3};
  % level 2: slot 1
  \node[n] (c1) at (-1.2,-2.6) {2};
  \node[n] (c2) at (0.2,-2.6) {3};
  \node[n] (c3) at (1.6,-2.6) {1};
  \node[n] (c4) at (3.0,-2.6) {3};
  \node[n] (c5) at (4.6,-2.6) {0};
  % level 3: slot 2
  \node[n] (d1) at (-1.2,-3.9) {1};
  \node[n] (d2) at (0.2,-3.9) {3};
  \node[n] (d3) at (3.6,-3.9) {0};
  \node[n,fill=black!8] (T) at (1.4,-5.2) {$\mathbf 1$};
  \draw[rt] (a1) -- (a2); \draw[rt] (a2) -- (a3);
  \draw[rt] (b1) -- (b2);
  \draw[rt] (c1) -- (c2); \draw[rt] (c3) -- (c4);
  \draw[rt] (d1) -- (d2);
  \draw[dn] (a1) -- (b1); \draw[dn] (a2) -- (b3); \draw[dn] (a3) -- (b4);
  \draw[dn] (b1) -- (c1); \draw[dn] (b2) -- (c1);
  \draw[dn] (b3) -- (c3); \draw[dn] (b4) -- (c5);
  \draw[dn] (c1) -- (d1); \draw[dn] (c2) -- (d1);
  \draw[dn] (c3) -- (d3); \draw[dn] (c4) -- (d3); \draw[dn] (c5) -- (d3);
  \draw[dn] (d1) -- (T); \draw[dn] (d2) -- (T); \draw[dn] (d3) -- (T);
  \node[font=\scriptsize,anchor=west] at (5.4,0) {shape};
  \node[font=\scriptsize,anchor=west] at (5.4,-1.3) {slot 0};
  \node[font=\scriptsize,anchor=west] at (5.4,-2.6) {slot 1};
  \node[font=\scriptsize,anchor=west] at (5.4,-3.9) {slot 2};
\end{tikzpicture}
\caption{The reachable states of \cref{ex:sym-par} as an LDD (derived by
hand; the five padding positions of the real vectors, each a single node
$0$, are left out). Solid: down; dashed: right. Values: $1 = \code{pre a
nil}$, $2 = \code{pre b nil}$, $3 = \code{nil}$, $0$ = no slot; shapes
$1 = \sigma_1$, $2 = \sigma_2$, $3 = \sigma_3$. The 15 nodes hold 11
vectors: the 8 of $\{1\}\times\{1,3\}\times\{2,3\}\times\{1,3\}$, then
$(2,3,1,0)$, $(2,3,3,0)$, $(3,3,0,0)$. The nodes $2{\to}3$ at level 2 and
$0$ at level 3 are shared.}
\label{fig:sym-ldd}
\end{figure}

\subsubsection{Relations}

A relation is an LDD read through a \emph{meta} vector, one code per state
position, ended by $-1$ (\code{0xFFFFFFFF};
\locr{lib/ldd/ldd.ml}{73}{84}). For a set $S$ of vectors of length $n$,
a relation LDD $R$ and meta $m$, Sylvan's \code{lddmc_relprod} and
\code{lddmc_relprev} interpret each position $i$ by its code:

\begin{center}
\small
\begin{tabular}{@{}c l l l@{}}
\toprule
code & \code{Ldd.level} & levels of $R$ & meaning for $(w, w')$ \\
\midrule
0 & \code{Copy} & none & $w'_i = w_i$ \\
1, 2 & \code{Read}, \code{Write} & two: old, new & $R$ holds $w_i$, then $w'_i$ \\
3 & \code{Only_read} & one & $R$ holds $w_i$, and $w'_i = w_i$ \\
4 & \code{Only_write} & one & $R$ holds $w'_i$, $w_i$ arbitrary \\
5 & \code{Label} & one, no state position & skipped \\
$-1$ & (end) & & the remaining positions copied \\
\bottomrule
\end{tabular}
\end{center}

Writing $w \to_{R,m} w'$ when some sequence of $R$ matches $(w,w')$
under $m$:
\[
\mathrm{relprod}(S,R,m) = \{w' \mid \exists w\in S.\ w\to_{R,m}w'\},
\]
\[
\mathrm{relprev}(S,R,m,U) = \{w \in U \mid \exists w'\in S.\ w\to_{R,m}w'\}.
\]
The point of the meta is locality: a relation over $k$ touched positions
of an $n$-vector stores only those $k$ (or $2k$) levels, and one learned
tuple applies to every state that agrees on the read positions.
\code{Label} is bound and tested (\loc{test/ldd_tests.ml}{47}) but the
engine never uses it: it keeps one relation per label instead.

\subsubsection{The bound operations}

Every OCaml function and the C stub behind it
(\file{lib/ldd/ldd.mli}; \file{lib/ldd/ldd.ml} lines 5--24 for the
externals; \file{lib/ldd/ldd_stubs.c}):

\begin{center}
\small
\begin{tabularx}{\linewidth}{@{}>{\raggedright\arraybackslash}p{3.2cm} >{\raggedright\arraybackslash}p{3.6cm} >{\raggedright\arraybackslash}X@{}}
\toprule
OCaml (\file{ldd.mli}) & C stub (\file{ldd_stubs.c}) & Sylvan call, notes \\
\midrule
\code{init} (22) & \code{ml_ldd_init} (74--84) & \code{lace_start(workers, 1000000)};
  \code{sylvan_set_limits(mb MB, 1, 5)}; \code{sylvan_init_package};
  \code{sylvan_init_ldd}. Only the first call acts (static \code{started}).
  Defaults: 1 worker, 1024\,MB (\loc{lib/ldd/ldd.ml}{26}). \\
\code{empty}, \code{epsilon} (25, 28) & \code{ml_ldd_false}, \code{_true} (86--87) &
  $\mathbf 0$, $\mathbf 1$. \\
\code{cube} (32) & \code{ml_ldd_cube} (89--95) & \code{lddmc_cube} inside a Lace task
  (\code{cube_task}, 56); values checked in OCaml to lie in $[0,2^{32}-2]$
  (\locr{lib/ldd/ldd.ml}{29}{40}); width at most 4096 (\code{MAX_WIDTH}, 20). \\
\code{of_list} (36) & \code{ml_ldd_union_cubes} (97--104) & Blocks of 64 cubes by
  \code{lddmc_union_cube} in a task (58--70), merged by halves with
  intermediate sets released (\locr{lib/ldd/ldd.ml}{42}{58}). Row width is
  taken from the first row and not checked for the others. \\
\code{union}, \code{minus}, \code{inter} (38--40) & 106--122 &
  \code{lddmc_union}, \code{lddmc_minus}, \code{lddmc_intersect}. \\
\code{equal}, \code{is_empty} (41--42) & \code{ml_ldd_equal} (146) & Node identity
  (canonicity); \code{is_empty s} is \code{equal s (empty ())}. \\
\code{mem} (43) & \code{ml_ldd_mem} (148--153) & \code{lddmc_member_cube}. \\
\code{count} (46) & \code{ml_ldd_count} (142) & \code{lddmc_satcount} (a long double)
  as a float. \\
\code{nodes} (49) & \code{ml_ldd_nodes} (144) & \code{lddmc_nodecount}. \\
\code{elements} (52) & \code{ml_ldd_elements} (188--207) & A read-only walk
  (\code{walk}, 168--186) collecting at most \code{limit} rows
  (default 1000, \loc{lib/ldd/ldd.ml}{61}), in lexicographic order; the
  engine always passes \code{~limit:max_int}. \\
\code{rel_meta} (65) & via \code{cube_c} & The meta cube, codes then $-1$. \\
\code{relprod}, \code{relprev} (68, 72) & 124--134 & \code{lddmc_relprod},
  \code{lddmc_relprev} (with \code{within} as Sylvan's universe). \\
\code{project} (76) & \code{ml_ldd_project} (136--140) & \code{lddmc_project} with
  codes 1 keep, 0 quantify, $-2$ ``quantify the rest''
  (\locr{lib/ldd/ldd.ml}{91}{100}). \\
\code{product} (80) & \code{ml_ldd_join} (209--213) & \code{lddmc_join} with
  projections $1^{w}\,{-2}$ for $a$ and $0^{w}\,{-1}$ for $b$: since $b$
  has none of $a$'s positions, the join is $\{u\cdot v\}$
  (\locr{lib/ldd/ldd.ml}{104}{115}). \\
\code{release} (83) & \code{ml_ldd_release} (155--162) & \code{lddmc_deref} now;
  the value then holds $\mathbf 0$. \\
\code{table_usage} (87) & 217--227 & \code{sylvan_table_usage}'s filled count, and
  \code{llmsset_get_max_size} (the maximum, not the current size). \\
\code{gc} (91) & \code{ml_ldd_gc} (230--235) & \code{sylvan_gc}. \\
\code{reclaim} (96) & OCaml only & If the table is over half full,
  \code{Gc.full_major ()} then \code{gc ()} (\locr{lib/ldd/ldd.ml}{117}{123}). \\
\bottomrule
\end{tabularx}
\end{center}

Sylvan's own operations (\code{lddmc_union} etc.) are macros that
\code{RUN} a Lace task (\file{vendor/sylvan/src/sylvan_ldd.h}, e.g.\
line 106), so the stubs may call them from the OCaml thread; only
\code{lddmc_cube} and \code{lddmc_union_cube} are wrapped by hand
(\locr{lib/ldd/ldd_stubs.c}{1}{8}).

\subsubsection{Memory: who keeps what alive, and when it is freed}

\begin{itemize}
\item \textbf{One reference per OCaml value.} \code{wrap}
  (\locr{lib/ldd/ldd_stubs.c}{40}{46}) allocates a custom block with
  \code{caml_alloc_custom_mem(.., 4096)}, telling OCaml's GC that each value
  holds about 4\,KB outside the heap (so the major GC runs sooner), and
  takes \code{lddmc_ref} on the node. The finaliser
  (\locr{lib/ldd/ldd_stubs.c}{22}{25}) drops the reference.
  \code{compare} and \code{hash} are by node index
  (\locr{lib/ldd/ldd_stubs.c}{27}{33}), so OCaml's \code{=} on \code{Ldd.t}
  is set equality.
\item \textbf{Nodes are freed only by Sylvan's GC}, which runs on demand
  when the unique table fills during an operation, or by \code{Ldd.gc}. It
  frees every node not reachable from a reference. A dead OCaml value not
  yet finalised still holds its reference, so its nodes stay. Hence
  \code{release} (free now) and \code{reclaim} (finalise, then collect).
\item \textbf{When the table is full even after a collection,} Sylvan
  prints ``MDD Unique table full'' and calls \code{exit(1)}
  (\file{vendor/sylvan/src/sylvan_ldd.c}, lines 373--374 and 400--401). Inside
  Rocq this ends the whole process (and the proof with it); it is not an
  exception, so no tactic can report it.
  The size is fixed at the first \code{init}: 1024\,MB in the plugin
  (\loc{lib/learn/learn.ml}{858} calls \code{Ldd.init ()}), 256\,MB in
  \file{test/ldd_tests.ml}. There is no Rocq option to change it.
\item \textbf{Sylvan is never shut down.} Once started inside a Rocq
  process, its tables (and one Lace worker thread) live until the process
  ends; nodes from one command are collected by later GCs once their OCaml
  values are finalised.
\end{itemize}

What the engine releases explicitly is listed with each algorithm below.
The rewrite of the game for memory (Phase~4) was prompted by three-process
coarse-grained against fine-grained CADP filling the 1024\,MB table in the game's sixth
iteration, while its relation was only 114k nodes: dead intermediate sets
filled the rest.

\subsubsection{Building and testing}

\file{lib/ldd/dune} copies the C files of \file{vendor/sylvan/src} and
\file{vendor/lace/src} and compiles them
\emph{into the binding's own library} (\code{-O2 -std=gnu11
-D_GNU_SOURCE -w}), not into a separate C library: Rocq loads the plugin's
libraries as \code{.cmxs} with Dynlink, and a C archive's objects reach a
\code{.cmxs} only when the same link refers to them.
\file{lib/ldd/README.md} gives the upstream versions, the two
hand-written configuration headers (\file{lib/ldd/config/sylvan_config.h},
\file{lace_config.h}) and how to update. \file{test/ldd_tests.ml} checks counts, the set
operations, canonicity, large values, image and preimage with a copy
level, a labelled relation, projection, product, release, and locality:
two independent counters mod 100, one relation each, reach 10,000 states
in at most 202 nodes (\locr{test/ldd_tests.ml}{64}{89}).

\subsection{Learning relations: \texttt{lib/learn}}\label{sec:sym-learn}

\subsubsection{States as vectors}

All tables are shared by the two systems of a check, so that labels
compare (\code{tables}, \locr{lib/learn/learn.ml}{40}{52}): values,
shapes and labels each get an integer from 1 in order of first meeting
(\code{intern}, \locr{lib/learn/learn.ml}{91}{99}); 0 means ``no slot''.
Terms are hashed by \code{deep_hash} (\locr{lib/learn/learn.ml}{8}{21}),
a structural hash of the whole term consistent with \code{Constr.equal}:
\code{Hashtbl.hash} looks at 10 nodes only, and CADP's large terms that
differ deep down collided (stage 2c).

A state $t$ of a system with \code{width} $w$ is the vector
\[
\mathrm{enc}(t) = [\,\#\mathit{shape};\ \#v_0;\ \dots;\ \#v_{k-1};\
0;\ \dots;\ 0\,] \in \N^{w+1}
\]
from \code{Skeleton.split} (\code{encode},
\locr{lib/learn/learn.ml}{117}{138}; \code{decode} rebuilds,
\locr{lib/learn/learn.ml}{140}{146}), each slot at its place in the
shape's order (below). If a state has more than $w$ slots,
\code{Too_wide n} is raised (also from \code{make_summand} and
\code{split_target}), and \code{check} restarts that system's exploration
with width $\max(n, 2w)$ (\code{run_explore},
\locr{lib/learn/learn.ml}{1632}{1635}). The initial width is 8.

\paragraph{The order of the slots.} A slot need not sit at its own
index: since 9 November 2026, for a system whose summands all keep the
start's shape, the slots go in order of how many summands read or write
them, most first, ties left to right (\code{explore},
\locr{lib/learn/learn.ml}{1141}{1187}). Each summand's positions are
then near the front and the rest of the vector is copied, which is what a
static variable ordering over the dependency matrix aims
for~\cite{meijer-vdpol16}; in CADP the shared memory comes first, each
process's own slots after. The order is a table per shape
(\code{perm}, \code{pos}; \code{pos_of} and \code{slot_at},
\locr{lib/learn/learn.ml}{69}{89}), read by \code{encode},
\code{decode}, the summands' codes and the partial terms, and handed to
the summands' codes and the partial terms. \code{MEBI_ORDER=off} keeps
each slot at its index.

\subsubsection{The interpreter}\label{sec:sym-interp}

A summand is learned on a \emph{partial term}: the source shape with the
read slots' values and, in every other slot $i$, a placeholder
\code{Meta i} (\code{partial}, \locr{lib/learn/learn.ml}{994}{1008}).
\code{fire c lts src chain} (\locr{lib/learn/learn.ml}{295}{325}) follows
the chain:

\begin{shell}[mathescape=true]
fire(lts, src, chain):                                   (learn.ml:259-312)
  []           -> []
  "<..>" :: r  -> if src has a placeholder: Stuck
                  else [(l, t, r) | (l, t) in steps(lts, src)]
  name :: r    -> $\rho$ := the rule called name;  b := cmatch($\rho$.src, src)   (No_match -> [])
                  for each b' in prems($\rho$.premises, b, r):
                      (compute(subst b' $\rho$.lbl), compute(subst b' $\rho$.tgt), rest)
prems(Decide w :: ps)        -> if compute(subst b w) = true then prems(ps) else []
prems(Opaque _ :: _)         -> Stuck
prems(Step L' s l t :: ps)   -> for each (l', t', r') in fire(L', compute(subst b s), chain):
                                   cmatch l against l', t against t'; prems(ps)
steps(lts, src):  memoised per (lts, src)
  if depth $\ge$ 64: Unfollowable(lts, src)
  else interp(lts, src), or Unfollowable(lts, src) on Stuck   (interpreted += 1)
interp(lts, src): every rule of lts, as fire, premises by prems_all  (learn.ml:345-385)
\end{shell}

\begin{itemize}
\item \code{cmatch} (\locr{lib/learn/learn.ml}{167}{202}) is first-order
  matching of a rule's source pattern against a term. A constructor
  pattern against a placeholder raises \code{Stuck}; a variable bound twice
  to a term with a placeholder raises \code{Stuck}; a closed non-constructor
  pattern (a constant such as \code{Acquire}) is compared with the term
  after \code{nf_all}.
\item \code{compute} (\locr{lib/learn/learn.ml}{212}{227}) normalises a
  closed term with \code{nf_all}; around placeholders it only descends
  through constructor applications, and raises \code{Stuck} on anything
  else or on \code{unbound} (\code{Meta 1000000}, a variable nothing
  bound, \loc{lib/learn/learn.ml}{199}). So a placeholder in the result
  means ``that slot carried unchanged''.
\item Two exceptions with distinct meanings: \code{No_match} (this rule
  does not apply) and \code{Stuck} (the interpreter cannot decide: on a
  closed term, the step is \code{Unfollowable}, an error). \code{max_depth} $= 64$ (\loc{lib/learn/learn.ml}{281}) stops a
  rule whose premise steps from its own source from recursing forever.
\item Rules of other LTSs are read lazily through \code{rules_of}
  (\locr{lib/learn/learn.ml}{283}{291}).
\end{itemize}



\subsubsection{Summands' relations and meta codes}

Each summand of a source shape becomes a record
(\locr{lib/learn/learn.ml}{410}{420}) with its chain, its projection
\code{keep} (position 0, the shape, always; then the read slots), its
static codes, the projections already learned (\code{seen}), its
relations, and the projections that escaped. The static codes come from
\code{make_summand} (\locr{lib/learn/learn.ml}{443}{500}). With $k_s$ the
source's slots, $k_t$ the target's, and \emph{same} whether the target
shape equals the source shape:
\[
\mathrm{reads} = \mathrm{reads}_{\mathrm{skeleton}} \cup \{ i \mid
  \text{target slot } j \text{ is } \mathsf{Moved}\ i,\ \neg(\mathit{same}
  \wedge i = j) \}
\]
(a value moved to another place counts as read: its new place holds the
old value), and position $p$ gets (LTSmin's read, write and copy
dependencies~\cite{meijer14})
\[
\mathrm{code}(0) = \begin{cases}\mathsf{RO} & \mathit{same}\\
\mathsf{RW} & \text{else}\end{cases}
\qquad
\mathrm{code}(1+j) = \begin{cases}
\mathsf{RW} & j\in\mathrm{reads},\ \mathrm{written}(j)\\
\mathsf{RO} & j\in\mathrm{reads},\ \neg\mathrm{written}(j)\\
\mathsf{WO} & j\notin\mathrm{reads},\ \mathrm{written}(j)\\
\mathsf{Copy} & \text{else}
\end{cases}
\]
where $\mathrm{written}(j)$ holds if $j<k_t$ and target slot $j$ is not
$\mathsf{Moved}\ j$ in the same shape, or $k_t \le j < k_s$ (the slot
disappears and is padded to 0). \code{levels}
(\locr{lib/learn/learn.ml}{469}{477}) maps $\mathsf{RW}$ to
\code{[Read; Write]}, $\mathsf{RO}$ to \code{Only_read}, $\mathsf{WO}$ to
\code{Only_write}, $\mathsf{Copy}$ to \code{Copy}.

\subsubsection{Learning a projection}

\code{learn sm frontier visited} (\locr{lib/learn/learn.ml}{625}{746}):

\begin{shell}[mathescape=true]
fresh := project(frontier, sm.keep) $\setminus$ sm.seen
for each pv in fresh with pv[0] = sm.sid:                   projections += 1
  results := fire(lts, partial(shape, sm.keep, pv), sm.chain)    (Stuck -> escape pv)
  for each (l, t): split t into (target shape, slots: Mv i | Val v)    (split_target, 523-535)
  if some Mv i with i $\ne$ j and slot i unread: escape pv     (a move of an unread slot)
  else for each result:
    widen the static codes where this result writes more        (needs, 695-716):
      RO $\to$ RW, Copy $\to$ WO   if position written;  WO $\to$ Copy  if not
    tuple := per code: RW (old,new), RO old, WO new, Copy nothing
    batch[(codes, label)] += tuple
sm.seen $\cup$= fresh;  add batch as one LDD per (codes, label)       (add_rel, 491-500)
escaped projections: learn their whole states among visited  (learn_global, 589-618)
\end{shell}

\begin{itemize}
\item A position ``needs'' a write (\locr{lib/learn/learn.ml}{695}{706})
  if it is the shape and the shape changed, a slot moved from elsewhere, a
  value that is not the read slot's own old value, or a slot beyond the
  new shape but within the old (padded to 0). A read slot comes back from
  the interpreter as a value, not a placeholder, so ``value equal to the
  old one'' is no write: this misclassification caused all 37,584 of
  Test4's escapes in stage 2b (\file{doc/learn.md}, stage 2c item 1).
\item Hence one summand may hold several relations, one per (codes,
  label): a result whose shape changed at run time (a dynamic summand,
  e.g.\ a \code{tfix} unfolding into a \code{tpar}) goes into a relation
  with more positions written.
\item \textbf{Escapes.} \code{learn_global} decodes each whole state and
  fires the chain on it; if that is \code{Stuck} too, it takes
  \emph{all} the state's steps from \code{steps} (a superset of the
  summand's, ``always correct to add''). Its tuples have every position
  $\mathsf{RW}$. Escapes cost time, never correctness: every tuple added is
  a real step. After stage 2c, escapes were 0 on every measured file.
\end{itemize}

\paragraph{A placeholder in a value is not checked.} \code{split_target}
interns any non-placeholder slot value as a value, even one that contains
a placeholder deeper inside (e.g.\ a non-structural constructor applied
to an unread slot). That would intern an open term as a state value.
Skeleton's \code{computed_reads} is meant to make every slot a computed
value depends on ``read'', so this should not arise; the code relies on
that analysis and does not check it (\locr{lib/learn/learn.ml}{526}{533}).

\begin{example}[\cref{ex:sym-par}, learned]
For \code{s_parl; s_parl; <step>} the codes are $[\mathsf{RO};
\mathsf{RW}; \mathsf{Copy};\dots]$ and the projections are $(1, v_0)$.
On $(1,1)$ the partial term is
\code{par (par (pre a nil) ?1) ?2}; \code{fire} matches \code{s_parl}
twice, then \code{steps step (pre a nil)} interprets \code{s_pre}: label
\code{Some a}, target \code{nil}. The result
\code{par (par nil ?1) ?2} splits to $[\mathsf{Val}\,3; \mathsf{Mv}\,1;
\mathsf{Mv}\,2]$, so the relation for label \code{Some a} is the single
tuple $(1;\,1,3)$: ``in shape 1, slot 0 goes from 1 to 3, the rest is
copied''. It applies to all four states $(1,1,\ast,\ast)$ at once.
For \code{s_parl; s_end} every position $0..3$ is $\mathsf{RW}$ (the shape
changes, slot 2 moves to slot 1, slot 2 is padded); its relation, label
\code{None}, is $\{(1,2,\,3,3,\,3,v,\,v,0) \mid v \in\{1,3\}\}$, learned on
the projections $(1,3,3,v)$; other projections do not match
\code{par nil nil} and give nothing.
\end{example}

\subsubsection{Chained exploration}

\code{explore} (\locr{lib/learn/learn.ml}{972}{1242}) is breadth-first with
LTSmin's~\cite{kant15} \emph{chaining}: within a round, each summand in turn is learned
on and applied to the round's frontier grown by what the summands before
it found.

\begin{shell}[mathescape=true]
visited := frontier := {enc(start)}
while frontier $\ne\emptyset$:                                      (763-809)
  reclaim()
  sms := summands of the shapes in frontier     (766-771)
  cur := frontier;  added := $\emptyset$
  for sm in sms:
    learn(sm, cur, visited)
    f := $\bigcup_{r\in sm.rels}$ relprod(cur, r.rel, meta(r.codes)) $\setminus$ visited
    visited, cur, added $\cup$= f;  learn escaped projections' states in f
  frontier := added
\end{shell}

Released explicitly: the projection, \code{fresh}, the old \code{seen},
\code{next}, \code{f}, and the previous \code{visited}, \code{cur},
\code{added}, \code{frontier}. Not released (left to finalisers): the meta
cubes built per application (\loc{lib/learn/learn.ml}{785}) and the
shape projection (\loc{lib/learn/learn.ml}{769}). Chaining pays where
summands feed each other (Proc's Test1: 3--6 rounds) and not on CADP's
fine-grained model (2 summands, 81 rounds).

\paragraph{A bug fixed on 2026-11-08: 1000 shapes.} The shapes of a
frontier were listed with \code{Ldd.elements} at its default \code{limit}
of 1000. A frontier with more than 1000 distinct shapes would silently have
skipped the summands of the others; their states would have stayed in
\code{visited} without successors, giving a smaller reachable set and
possibly a wrong verdict (no measured example has more than 9 shapes).
Found by the first draft of this report; now every call in the file passes
\code{~limit:max_int} (\loc{lib/learn/learn.ml}{769}; also 730, 745, 759,
1087, 1104, 1118). \code{Ldd.elements}'s default is still 1000, a trap for
a future caller.

\paragraph{Why it is complete.} Each state enters some \code{cur} exactly
once, when it is new; every projection of it is then learned (or escaped
and learned whole), and a projection's results, once learned, are all its
steps. So a later state with an already-seen projection gets its
successors from the existing tuples, and no earlier state can gain
successors later. Soundness rests on every tuple being a real step: from
the interpreter, which follows the rules.

\paragraph{Saturation, measured and set aside}~\cite{ciardo01} (Phase~3; \file{doc/future-work.md},
``Saturation for the symbolic route''). On three-process coarse-grained CADP
against \code{spec_lts}: 119\,s interpreting summands (1.92M projections),
16\,s other learning, \textbf{0.25\,s applying relations}, 0.2\,s the
game. Saturation speeds up the application of relations, which is under
1\% of the time on every example.

\subsection{The game, symbolically: weak and strong}\label{sec:sym-game}

Let $\to_A$ and $\to_B$ be the learned relations (the union over
summands, per label), $\Reach_A$, $\Reach_B$ the explored sets, and
\[
\Uni = \Reach_A \times \Reach_B
\]
built by \code{Ldd.product} (\loc{lib/learn/learn.ml}{887}); a pair is the
concatenation of $A$'s vector and $B$'s. A relation of $B$ acts on pairs
with $|A|$ copy codes in front (\code{rels_of ~offset},
\locr{lib/learn/learn.ml}{890}{909}); a relation of $A$ needs no padding,
since positions after the meta's end are copied.

\paragraph{The silent label, and \code{?strong}.} \code{tau}
(\locr{lib/learn/learn.ml}{1327}{1343}) is an \code{int option}:
\begin{ocaml}
  let tau =
    let none = Rocqlib.lib_ref "core.option.None" in
    if strong
    then None
    else
      CTbl.fold (fun t i acc -> match Constr.kind (fst (Constr.decompose_app t)) with
        | Constr.Construct (c, _) when (* c is None *) ... -> Some i
        | _ -> acc) tb.labels None
\end{ocaml}
In weak mode ($\mathit{strong}=\mathsf{false}$, the default) it is the
intern number of the label whose head constructor is \code{None}, found by
scanning the shared label table, or \code{None} if no reachable step is
silent. In strong mode it
is \code{None}: no label is silent, and a \code{None}-labelled step is an
ordinary label like any other. Write $\tau$ for the silent label when
there is one; everything below is defined for both modes.

\begin{definition}[The game]
For $Z\subseteq\Uni$ and a side $X\in\{A,B\}$ (the other side's component
fixed):
\begin{align*}
\pre^X_l(Z) &= \{ (p,q)\in\Uni \mid X\text{'s component steps by } l
  \text{ into a pair of } Z\} \\
T_X(Z) &= \begin{cases} \mu Y.\ (Z\cap\Uni) \cup \pre^X_\tau(Y) &
  \text{weak, with a silent label } \tau\\
  Z\cap\Uni & \text{strong (or no silent label)} \end{cases} \\
\Ok^X_l(Z) &= \begin{cases} T_X(Z) & l = \tau \\
  T_X(\pre^X_l(T_X(Z))) & l \ne \tau \end{cases}\\
\Bad_{A\to B}(Z) &= \textstyle\bigcup_{l\in L_A} \pre^A_l(\Uni \setminus \Ok^B_l(Z))
\end{align*}
($T_X(Z)$: $X$ reaches $Z$ by $\tau^\ast$; $\Ok^X_l(Z)$: $X$ answers $l$
weakly into $Z$; in strong mode there is no $l=\tau$ case), and
$\Bad_{B\to A}$ symmetrically. With $F_{\mathrm{sim}}(Z) = Z\setminus
\Bad_{A\to B}(Z)$ and $F_{\mathrm{bis}}(Z) = Z \setminus (\Bad_{A\to
B}(Z)\cup\Bad_{B\to A}(Z))$, the result is $Z^\ast = \nu F$, computed as
$Z_0 = \Uni$, $Z_{i+1} = F(Z_i)$ until $Z_{i+1} = Z_i$. The verdict is
$(\mathrm{enc}(s_A),\mathrm{enc}(s_B)) \in Z^\ast$.
\end{definition}

Unfolded: $(p,q)\in\Bad_{A\to B}(Z)$ iff $p\to^l_A p'$ for some $l, p'$
and there is no $q'$ with $q \wtr{l} q'$ (for $l=\tau$: $q\to_\tau^\ast
q'$) and $(p',q')\in Z$. This is exactly \code{simF}
(\locr{theories/Bisimilarity.v}{60}{64}): a strong move answered by a
\code{weak} move (\locr{theories/Bisimilarity.v}{44}{47}), the silent
case allowing zero steps; \code{bisimF}
(\locr{theories/Bisimilarity.v}{212}{218}) adds the converse.

\paragraph{The strong game.} With no silent label, $T_X(Z) = Z\cap\Uni$, so
$\Ok^X_l(Z) = \pre^X_l(Z)\cap\Uni$ for every $l$, and
\[
F^{s}_{\mathrm{sim}}(Z) = Z \setminus \textstyle\bigcup_{l\in L_A}
  \pre^A_l\bigl(\Uni\setminus\pre^B_l(Z)\bigr),
\qquad
F^{s}_{\mathrm{bis}}(Z) = F^{s}_{\mathrm{sim}}(Z) \setminus
  \textstyle\bigcup_{l\in L_B}\pre^B_l\bigl(\Uni\setminus\pre^A_l(Z)\bigr).
\]
Unfolded: $(p,q)$ is removed iff $p\to^l_A p'$ and no $q\to^l_B q'$ has
$(p',q')\in Z$ (and, for bisimilarity, symmetrically). The result is again
$Z^\ast = \nu F^{s}$, iterated from $\Uni$. Its post-fixpoints are exactly
the relations satisfying \code{ssimF} (\code{sbisimF}) of
\locr{theories/Bisimilarity.v}{361}{376}: a step, \code{None} included,
answered by one step with the same label. Since the same code computes both,
the strong game is the weak one with $\tau^\ast$ replaced by the identity;
nothing else changes (\code{bad} still takes the answerer's
\code{tau_star} once, which is now $Z\cap\Uni$, and the test
\code{Some l = tau} is never true).

\begin{proposition}
$F$ is monotone and $F(Z)\subseteq Z$, so the iteration from $\Uni$
reaches $\nu F$ in at most $|\Uni|$ steps. $\nu F$ is the largest weak
simulation (for $F_{\mathrm{bis}}$: the largest weak bisimulation)
contained in $\Uni$ of the learned LTSs. If the learned relations are
exactly the Rocq LTS's steps on reachable states, the verdict is
\code{weak_sim} (\code{weak_bisimilar}) of the start terms; for
$F^{s}$, \code{strong_sim} (\code{strong_bisimilar}).
\end{proposition}
\begin{proof}[Sketch]
$Z\subseteq Z'$ gives $T_X(Z)\subseteq T_X(Z')$, hence
$\Ok(Z)\subseteq\Ok(Z')$ and $\Bad(Z)\supseteq\Bad(Z')$. A post-fixpoint
$R\subseteq F(R)$ is a relation in $\Uni$ whose every move is answered
into $R$, i.e.\ a weak (bi)simulation; the greatest fixpoint contains all
of them. Every successor of a reachable state is reachable, so
restricting to $\Uni$ loses no pair a (bi)simulation from the start pair
needs.
\end{proof}
This argument is not mechanised. What is checked in Rocq is the
certificate (\cref{sec:prooflog}), not the engine.

\paragraph{With LDD operations} (\locr{lib/learn/learn.ml}{927}{1024}):
\begin{shell}[mathescape=true]
pre(X, l, S)       = $\bigcup$ over X's relations r with label l: relprev(S, r, meta, within = U)
tau_star(X, Z)     = Z := Z $\cap$ U;  if tau = None: Z            -- strong: the identity on U
                     fresh := Z
                     while fresh $\ne\emptyset$: p := pre(X,$\tau$,fresh); fresh := p $\setminus$ Z; Z $\cup$= fresh
bad(M, Ans, Z)     = tz := tau_star(Ans, Z)                 -- once per iteration, not per label
                     $\bigcup_{l\in L_M}$ pre(M, l, U $\setminus$ ok_l),   ok_$\tau$ = tz,  ok_l = tau_star(Ans, pre(Ans, l, tz))
                     -- strong: tz = Z $\cap$ U and ok_l = pre(Ans, l, tz) $\cap$ U for every l
gfp(Z)             = bd := bad(A, B, Z) [$\cup$ bad(B, A, Z) for Bisim]
                     Z' := Z $\setminus$ bd;  if Z' = Z then Z else gfp(Z')
verdict            = mem(gfp(U), start)
\end{shell}
A pair is one vector, the \emph{answering} side's positions first (for
\code{Sim}, $B$'s: $\Uni = \Reach_B \times \Reach_A$; for \code{Bisim},
where both answer, the narrower side's; \code{b_first},
\loc{lib/learn/learn.ml}{1305}), so that the answerer's relations, applied
most, need no copied prefix; \code{start} is the two start vectors in that
order. A relation applied at an offset copies the positions before it, and
its meta stops at its last touched position (\code{meta_levels},
\loc{lib/learn/learn.ml}{520}): past the end Sylvan carries the rest of
the vector unchanged.
\code{tau_star} (\locr{lib/learn/learn.ml}{1354}{1373}; the strong case is
\loc{lib/learn/learn.ml}{1356}) applies only the pairs new in each round.
\code{pre} is at \locr{lib/learn/learn.ml}{1344}{1350}, \code{labels_of} at
\loc{lib/learn/learn.ml}{1375}, \code{bad} at
\loc{lib/learn/learn.ml}{1385}, \code{gfp} at
\loc{lib/learn/learn.ml}{1414}, called from $\Uni$ at
\loc{lib/learn/learn.ml}{1444}; the verdict is \code{Ldd.mem} of the start
pair (\locr{lib/learn/learn.ml}{1555}{1560}). Each iteration starts with
\code{reclaim ()}. Intermediate sets are released as soon as they are
dead (\code{union_into}, \loc{lib/learn/learn.ml}{1337}); the meta
cubes in \code{rels_tbl} and \code{uni} are not. Each iteration recomputes
the whole of $\Bad(Z)$; nothing is incremental across iterations
(40--47 iterations on CADP).

\subsection{\texttt{Learn.check} and its report}\label{sec:sym-check}

\begin{ocaml}
val check : Environ.env -> Evd.evar_map -> ?width:int
  -> ?strong:bool -> kind -> system -> system -> report
val pr : ?strong:bool -> kind -> report -> Pp.t
\end{ocaml}
(\file{lib/learn/learn.mli}.)
\code{kind} is \code{Bisim} or \code{Sim} (the first simulated by the
second); \code{strong} (default \code{false}) selects the strong game
(\cref{sec:sym-game}); it changes nothing before the game, so exploration,
the report's counts and the reachable sets are the same in both modes.
The tactics pass \code{~strong:true} for \code{sim}/\code{bisim} and
\code{false} for \code{weak_sim}/\code{weak_bisim}. \code{check} starts
Sylvan, builds the shared tables and interpreter context, explores $A$
then $B$, builds $\Uni$, runs the game, and, with the log on
(\code{proof_log}), repeats the last check with Sylvan's proof log open
(\cref{sec:prooflog}). The \code{report}:

\begin{center}
\small
\begin{tabularx}{\linewidth}{@{}l >{\raggedright\arraybackslash}X@{}}
\toprule
field & meaning \\
\midrule
\code{verdict} & start pair $\in Z^\ast$ \\
\code{states} & $|\Reach_A|$, $|\Reach_B|$ (floats, from \code{count}) \\
\code{pairs}, \code{nodes} & $|Z^\ast|$ and its node count \\
\code{labels} & distinct labels, both systems \\
\code{interpreted} & closed terms interpreted (shared context; counts
  from an aborted \code{Too_wide} attempt are included) \\
\code{projections}, \code{escapes} & per system \\
\code{summands}, \code{width}, \code{rounds} & per system; \code{width} is
  the vector length $w+1$ \\
\code{iterations}, \code{seconds} & of the game; wall-clock of the whole
  check \\
\bottomrule
\end{tabularx}
\end{center}

Exceptions: \code{Unfollowable}; and, outside OCaml, Sylvan's
\code{exit(1)} on a full table.

\subsection{The \texttt{debug} output}\label{sec:sym-debug}

Under \code{debug} $t$ for any of the four tactics, \code{prove} runs the
same \code{Learn.check} and, before testing the
verdict, prints \code{debug }$t$\code{: } followed by \code{Learn.pr
~strong k r} (\locr{lib/learn/learn.ml}{1142}{1195}) as a notice
(\file{src/user_tactics.ml}). The tactic then proceeds exactly as
without \code{debug}. There is no ``not run'' case any more: an exception
from the engine is the tactic's error. The format:
\begin{shell}
debug weak_bisim: symbolic (learned per summand, Sylvan): weakly bisimilar.
  reachable states: <a> and <b>; pairs in the relation: <n> (<k> nodes); labels: <l>
  summands: <sa> and <sb>; projections learned: <pa> and <pb> (escaped: <ea> and <eb>); steps interpreted: <i>
  vector widths <wa> and <wb>; exploration rounds: <ra> and <rb>; game iterations: <g>; <t>s
\end{shell}
The word is ``strongly'' under \code{~strong:true}
(\loc{lib/learn/learn.ml}{1143}); the verdict reads ``NOT \ldots\
bisimilar'', ``\ldots\ simulated'' or ``NOT \ldots\ simulated'' otherwise
. The goal checks of \code{prove}
(closed terms, LTS-shaped inductives) run first, so the engine never sees
an open goal.

\subsection{Limits}\label{sec:sym-limits}

\begin{itemize}
\item A closed premise's steps are interpreted from its LTS's rules to a
  depth of 64; a step the interpreter cannot follow (a rule whose source
  is not a pattern it can match, or a premise stepping from its own source
  without end) is an error, \code{Learn.Unfollowable}.
\item An LTS must be an inductive \code{M -> option A -> M -> Prop} without
  parameters.
\end{itemize}

\endgroup
